%% sections/06-discussion.tex

\section{Discussion}
\label{sec:discussion}

\subsection{The Pattern as Pipeline Validation}

The NPC arc-completion pattern is, in part, a quality signal for the adventure design
pipeline. When NPCs have designed arc-completions, they fire by character logic. When
they do not, the DM must improvise, and improvised moments are weaker. The 90\% session
coverage rate is not primarily evidence that the pattern works at runtime; it is evidence
that the NPC-architect pipeline is functioning at design time.

This reframes the pattern's design implication: the primary intervention is not
``tell DMs to recognize and time NPC moments better,'' but ``tell adventure designers
to specify NPC arc-conditions before play.'' The constraint produces reliability
because it forces the designer to articulate what the NPC is actually capable of saying,
under what conditions, and what the DM must not do to compromise the moment.

\subsection{Why Small Moments Are Better}

The small-moment principle has a structural explanation in terms of narrative pacing.
When a party has accumulated 3--5 sessions of investment in an NPC relationship,
the moment of arc-completion is the release of that accumulated tension. A compressed
arc-completion allows the tension to release completely --- one line, one act, space.
An extended arc-completion redistributes the tension back into exposition, diluting
the release.

This pattern is familiar from literary and dramatic analysis~\cite{chekhov1886letter,
aristotle_poetics}: the most emotionally powerful moments in well-crafted narratives
are often the briefest. What makes them powerful is the accumulated context, not the
moment itself. The moment needs only to be the hinge that turns everything that came
before into meaning.

Design implication: when writing arc-completion text, the designer should aim for the
shortest formulation that is semantically complete. If the arc-completion text is more
than two sentences, it is probably not the arc-completion; it is the DM's wish for
what should happen afterward.

\subsection{Morally Grey NPCs: Design Considerations}

The 31\% morally-grey proportion in arc-completion instances suggests that the pattern
is not limited to sympathetic NPCs. Gorret (Campaign 3, first-session antagonist
faction member), Aldras Bent (Campaign 5, guild fence with divided loyalties), and
Saren Coldforge (Campaign 4, military antagonist) all fired arc-completions that
were among the most scored moments in their respective campaigns.

The key design principle for morally grey arc-completions is that the arc-completion
text should be \textit{consistent with the NPC's moral valence}, not a conversion from it.
Gorret does not become friendly when he says ``at least you say what you are.''
He acknowledges something. His moral position is unchanged; the acknowledgment is
within it. Saren's silent token-acceptance is not a surrender; it is a recognition
that something she has been carrying for twenty years can finally be set down.

NPCs who undergo moral conversion in their arc-completion tend to produce less
satisfying moments than NPCs who express their existing moral position with greater
clarity. The arc-completion is a moment of self-expression, not character transformation.

\subsection{Limitations and Future Work}

The corpus is AI-simulated: the DM has no capacity for improvisation and therefore
no temptation to script moments prematurely. In human-run TRPG, the pattern requires
DM discipline to hold: not extending the moment, not delivering the arc-completion
before the firing condition is met, not replacing a designed arc-completion with an
improvised one. We cannot report data on whether human DMs maintain this discipline.
Future work should apply the pattern in human-run sessions and measure compliance
with the DM discipline note.

Additionally, the corpus covers only D\&D 5e mechanics. Arc-completions in Powered
by the Apocalypse games (where NPC moves are explicitly designed and triggered) may
show different patterns. The pattern's core mechanism --- specify the condition and
the text before play; let character logic fire it --- should transfer, but the
specific design template may require adaptation.
